\section{Introdução}

Algoritmos Genéticos (AGs) são heurísticas evolutivas baseadas em seleção natural, amplamente utilizadas para otimização em espaços de busca complexos~\cite{eiben2015introduction}. A capacidade de combinar exploração estocástica com operadores simples de reprodução torna esses algoritmos versáteis em aplicações que vão de problemas combinatórios a tarefas de engenharia e inteligência artificial.

Entre os elementos que influenciam o desempenho dos AGs, o \textit{controle populacional} — isto é, o modo como indivíduos são mantidos, removidos ou substituídos ao longo das gerações — desempenha papel central. Estratégias como substituição total, substituição parcial, elitismo e abordagens baseadas em ciclos de vida têm sido utilizadas na prática, mas a literatura apresenta resultados heterogêneos, conduzidos em contextos experimentais distintos e com ausência de padronização. Como consequência, ainda é difícil comparar, de forma direta e controlada, o impacto real de cada mecanismo no comportamento do algoritmo.

Essa lacuna motiva o presente estudo, cujo objetivo é \textbf{comparar sistematicamente diferentes estratégias de controle populacional em um mesmo cenário experimental}, mantendo constantes todos os demais componentes do AG (representação, seleção, cruzamento e mutação). Para isso, utiliza-se o problema das \textit{n}-rainhas como ambiente de teste, um benchmark clássico pela sua estrutura simples, escalabilidade controlada e ampla adoção em estudos de otimização.

A contribuição deste trabalho é dupla: (i) oferecer uma avaliação empírica rigorosa de quatro mecanismos de controle populacional — substituição total, substituição parcial, elitismo e uma abordagem baseada em idade — sob condições experimentais uniformes; e (ii) analisar como cada estratégia afeta taxa de sucesso, esforço evolutivo e custo computacional à medida que a dificuldade do problema aumenta.

Os experimentos são conduzidos com múltiplas execuções independentes para cada configuração, permitindo estimar médias, variâncias e intervalos de confiança, além de aplicar testes estatísticos para comparação entre estratégias. Essa estrutura experimental fornece evidências quantitativas que auxiliam na escolha de mecanismos mais adequados para AGs aplicados a problemas combinatórios semelhantes.
